\chapter{The Law of Sines and the Law of Cosines}
\label{ch:ch12-laws-sines-cosines}

\section{Beyond the Right Triangle}

The trigonometry developed so far has rested on two geometric settings, and
both, on inspection, share a common feature. The right triangle of
Chapter~3 supplied sine, cosine, and tangent as ratios invariant under
scaling, with a right angle present explicitly; the unit circle of
Chapter~4 extended those same functions beyond acute angles by measuring
coordinates against a pair of perpendicular axes, with the right angle now
hidden inside the coordinate system rather than inside a triangle. As long
as a right angle is available somewhere in the configuration, the tools
built in earlier chapters suffice. Many geometric situations, however,
supply no such angle. A surveyor measuring the distance between two
inaccessible points may obtain a triangle with no right angle at all; a
navigator determining position from bearings, or an astronomer inferring
distance from parallax, may encounter a triangle whose angles bear no
particular relationship to ninety degrees. If trigonometry is to serve
these settings, it must be extended beyond the right triangle, and this
chapter's two central results, the Law of Sines and the Law of Cosines,
should reduce to the familiar right-triangle relationships whenever a
right angle happens to be present, so that right-triangle trigonometry
becomes a special case of a broader theory rather than a separate subject
entirely.

\section{Deriving the Law of Sines}

Consider a triangle with sides $a$, $b$, $c$ opposite angles $A$, $B$,
$C$ respectively. Dropping an altitude of length $h$ from the vertex
opposite side $c$ divides the triangle into two right triangles sharing
that altitude, which provides a bridge between two different
trigonometric descriptions of the same quantity.

\begin{theorem}[Law of Sines]
For any triangle, $\dfrac{a}{\sin A} = \dfrac{b}{\sin B} = \dfrac{c}{\sin
C}$.
\end{theorem}

\begin{proof}
In the right triangle containing angle $A$, $\sin A = h/b$, so $h =
b\sin A$; in the right triangle containing angle $B$, $\sin B = h/a$, so
$h = a\sin B$. Both expressions describe the same altitude, so $a\sin B =
b\sin A$, and dividing by $\sin A \sin B$ gives $a/\sin A = b/\sin B$.
Repeating the construction with an altitude dropped from a different
vertex relates the third side-angle pair in the same way, establishing all
three ratios as equal.
\end{proof}

Each ratio in the theorem measures the same underlying quantity, namely
twice the circumradius of the triangle, so the equality of the three
ratios expresses a single consistency constraint that any valid
configuration of sides and angles must satisfy, rather than three
independent facts. The formula is also consistent with what Chapter~3
already established: if $C = 90^\circ$, then $\sin C = 1$ and the Law of
Sines reduces to $a/\sin A = b/\sin B = c$, or equivalently $\sin A = a/c$
and $\sin B = b/c$, exactly the right-triangle definitions of Chapter~3.
The Law of Sines therefore extends the earlier theory rather than
replacing it.

\begin{example}
Solve the triangle with $A = 42^\circ$, $B = 68^\circ$, $a = 12$. Since the
angles of a triangle sum to $180^\circ$, $C = 180^\circ - 42^\circ -
68^\circ = 70^\circ$. The Law of Sines gives
\[
\frac{12}{\sin42^\circ} = \frac{b}{\sin68^\circ}, \qquad
b = 12\,\frac{\sin68^\circ}{\sin42^\circ} \approx 16.6,
\]
and similarly $c = 12\,\sin70^\circ/\sin42^\circ \approx 16.8$, completing
the triangle without any altitude or right angle ever having been given
directly.
\end{example}

\section{The Ambiguous Case}

Most triangle-solving configurations determine either one triangle or
none. The Law of Sines admits a notable exception when two sides and a
non-included angle are known, traditionally called the SSA case. Fixing
one side and one angle and allowing the remaining given side to pivot
about its attached endpoint like a hinged rod shows why: depending on its
length relative to the altitude and the adjacent side, the free endpoint
may fail to reach the base at all, may reach it in exactly one position,
or may reach it in two distinct positions, yielding no triangle, one
triangle, or two triangles respectively. The ambiguity has an algebraic
source as well as a geometric one: it arises because a single sine value
corresponds to two different angles within one revolution, precisely the
phenomenon already encountered in solving equations of the form
$\sin\theta = 1/2$ in Chapter~11. The SSA case is therefore not an
isolated irregularity but a direct consequence of the geometry of the
sine function.

\begin{algorithmproc}{How to Determine the Number of Triangles in the SSA Case}
Given angle $A$ and sides $a$ and $b$, compute the altitude $h = b\sin A$
from the vertex opposite $b$. If $a < h$, no triangle exists, since side
$a$ is too short to reach the base. If $a = h$, exactly one triangle
exists, a right triangle with the right angle opposite $a$. If $h < a <
b$, two triangles exist, corresponding to the two angles $B$ and $180^\circ
- B$ satisfying $\sin B = b\sin A / a$. If $a \geq b$, exactly one triangle
exists, since the supplementary angle $180^\circ - B$ would make the angle
sum of the triangle exceed $180^\circ$.
\end{algorithmproc}

\begin{example}
Solve the triangle with $A = 30^\circ$, $a = 10$, $b = 18$. The Law of
Sines gives $10/\sin30^\circ = 18/\sin B$, and since $\sin30^\circ =
1/2$, this becomes $20 = 18/\sin B$, so $\sin B = 18/20 = 0.9$ and $B
\approx 64.2^\circ$. Because $h = b\sin A = 9$ and $a = 10$ satisfies $h <
a < b$, a second solution is expected: $B = 180^\circ - 64.2^\circ =
115.8^\circ$ also satisfies $\sin B = 0.9$, and both angles are compatible
with the given data. The first possibility gives $C = 180^\circ - 30^\circ
- 64.2^\circ = 85.8^\circ$, and the second gives $C = 180^\circ - 30^\circ
- 115.8^\circ = 34.2^\circ$; both triangles are valid, each containing a
side of length $10$, a side of length $18$, and an angle of $30^\circ$,
with the given data alone insufficient to distinguish between them.
\end{example}

Whenever the Law of Sines produces an angle through an inverse sine
computation, the supplementary angle should be checked against the
algorithm above before it is discarded, since in the SSA case it will
sometimes, though not always, correspond to a second valid triangle.

\section{Toward the Law of Cosines}

The Law of Sines cannot resolve every configuration. If all three sides
of a triangle are known and an angle is sought, the Law of Sines offers no
starting point, since it requires an already-known side-angle pair. A
different relationship is needed, and it is obtained by returning to
coordinate geometry, the strategy Chapter~2 introduced for translating
geometric questions into algebra.

\section{Deriving the Law of Cosines}

Place a triangle with sides $a$, $b$, $c$ and opposite angles $A$, $B$,
$C$ so that the vertex at angle $C$ lies at the origin and the side of
length $b$ extends along the positive $x$-axis to the point $(b,0)$. The
third vertex lies at distance $a$ from the origin, at angle $C$ from the
positive $x$-axis, so by the unit-circle coordinate formulas of Chapter~4
it has coordinates $(a\cos C, a\sin C)$.

\begin{theorem}[Law of Cosines]
For any triangle, $c^2 = a^2 + b^2 - 2ab\cos C$, and similarly $a^2 = b^2
+ c^2 - 2bc\cos A$ and $b^2 = a^2 + c^2 - 2ac\cos B$.
\end{theorem}

\begin{proof}
The side of length $c$ is the distance between $(b,0)$ and $(a\cos C,
a\sin C)$, so by the distance formula of Chapter~2,
\[
c^2 = (b - a\cos C)^2 + (a\sin C)^2 = b^2 - 2ab\cos C + a^2\cos^2 C +
a^2\sin^2 C.
\]
By the Pythagorean identity of Chapter~9, $a^2\cos^2 C + a^2\sin^2 C =
a^2$, so $c^2 = a^2 + b^2 - 2ab\cos C$. The remaining two forms follow by
relabeling the triangle's vertices and repeating the same argument.
\end{proof}

The three ingredients of this derivation are drawn from earlier chapters
of the book: the coordinate placement from Chapter~2, the trigonometric
coordinates from Chapter~4, and the Pythagorean identity from Chapter~9.
Setting $C = 90^\circ$ gives $\cos C = 0$, so the correction term vanishes
and $c^2 = a^2 + b^2$, the ordinary Pythagorean theorem; the Law of
Cosines therefore generalizes the Pythagorean theorem rather than
replacing it, with the familiar right-triangle relationship recovered
exactly when the included angle is a right angle.

\section{Solving Triangles with the Law of Cosines}

The most direct use of the Law of Cosines occurs when two sides and the
included angle are known, the SAS case.

\begin{example}
Solve for $c$ given $a = 8$, $b = 11$, $C = 50^\circ$. The Law of Cosines
gives
\[
c^2 = 8^2 + 11^2 - 2(8)(11)\cos50^\circ = 64 + 121 - 176(0.6428) \approx
71.9,
\]
so $c \approx 8.48$. Once all three sides are known, either the Law of
Sines or the Law of Cosines can determine the two remaining angles, since
the Law of Sines requires only one known side-angle pair to proceed and
one, $C$, is already known.
\end{example}

\begin{practicenow}
Solve for $c$ given $a = 13$, $b = 9$, $C = 120^\circ$.
\end{practicenow}

A second application arises when all three sides are known and an angle
is sought, the SSS case.

\begin{example}
Given $a = 7$, $b = 9$, $c = 11$, find angle $C$. Substituting into
$c^2 = a^2 + b^2 - 2ab\cos C$,
\[
121 = 49 + 81 - 126\cos C, \qquad 126\cos C = 9, \qquad \cos C =
\frac{9}{126} = \frac{1}{14},
\]
so $C = \arccos(1/14) \approx 85.9^\circ$. Unlike the SSA configuration of
Section~12.3, no ambiguity arises here: three side lengths satisfying the
triangle inequality determine a unique triangle, since $\arccos$ returns
exactly one angle in $[0,\pi]$ for any given ratio, with no supplementary
angle to check.
\end{example}

\begin{algorithmproc}{How to Solve an Oblique Triangle}
Given the known measurements of a triangle, determine which case applies.
If two angles and any side are known (AAS or ASA), find the third angle
from the angle sum and apply the Law of Sines directly. If two sides and
the included angle are known (SAS), apply the Law of Cosines to find the
third side, then use the Law of Sines or a second application of the Law
of Cosines to find the remaining angles. If all three sides are known
(SSS), apply the Law of Cosines to find one angle at a time. If two sides
and a non-included angle are known (SSA), apply the Law of Sines and use
the algorithm of Section~12.3 to determine whether zero, one, or two
triangles satisfy the data.
\end{algorithmproc}

\section{The Cosine Correction Term}

The formula $c^2 = a^2 + b^2 - 2ab\cos C$ is often committed to memory
without attention to its geometric content. The term $a^2 + b^2$ is
exactly what the Pythagorean theorem would predict if sides $a$ and $b$
met at a right angle, and the term $-2ab\cos C$ measures the departure of
the actual triangle from that special case. When $C = 90^\circ$, $\cos C =
0$ and the correction vanishes entirely; when $C < 90^\circ$, $\cos C$ is
positive and the correction subtracts from $a^2+b^2$, so the side opposite
an acute included angle is shorter than the right-triangle prediction;
when $C > 90^\circ$, $\cos C$ is negative and subtracting a negative
quantity lengthens the opposite side beyond that prediction. The Law of
Cosines can therefore be read as the Pythagorean theorem together with an
explicit correction recording how far the included angle departs from a
right angle.

\section{Area from Trigonometry}

Returning to the altitude construction that opened this chapter yields a
further consequence. With sides $a$ and $b$ enclosing angle $C$, and $b$
taken as the base, the altitude from the opposite vertex is $h = a\sin C$,
so the ordinary area formula $\text{Area} = \tfrac12(\text{base})(\text{height})$
becomes
\[
\boxed{\text{Area} = \frac12 ab\sin C},
\]
with the analogous formulas $\text{Area} = \tfrac12 bc\sin A$ and
$\text{Area} = \tfrac12 ac\sin B$ following by relabeling. Each of these
computes an area without requiring the altitude to be found explicitly,
since the sine function already encodes the necessary height information.

\begin{example}
Find the area of the triangle with $a = 12$, $b = 15$, $C = 40^\circ$.
\[
\text{Area} = \frac12(12)(15)\sin40^\circ \approx 90(0.6428) \approx 57.9
\]
square units, obtained without ever computing an altitude directly.
\end{example}

\section{Applications}

The two laws developed in this chapter recover distances and angles that
cannot be measured directly. Surveyors use triangulation to determine the
widths of rivers, the positions of landmarks, and the boundaries of large
parcels of land; navigators infer position from measured bearings and
distances; engineers analyze forces acting along non-perpendicular
directions; astronomers estimate the geometry of distant objects from
measurements taken at different locations. In each case certain distances
or angles are inaccessible to direct measurement, yet sufficient
constraints can be gathered to reconstruct them indirectly, and the
process is fundamentally reconstructive in the sense introduced in this
book's preface: measurements are not merely recorded but used to infer a
structure that was not itself observed.

\section{Looking Ahead}

The Law of Sines generalized the relationship between a triangle's sides
and its opposite angles; the Law of Cosines generalized the Pythagorean
theorem. Together the two provide a complete method for solving an
arbitrary triangle from any sufficient set of given measurements. The next
chapter moves from the geometry of individual triangles toward a broader
language for describing space, returning to the vectors of Chapter~2 with
new tools available. Vectors allow magnitude and direction to be treated
as a single object, and among their properties is an operation, the dot
product, whose defining formula will turn out to have exactly the
structure of the Law of Cosines derived here: a length-squared relation
with a cosine correction term. What in this chapter is a statement about
the sides and angles of a triangle will reappear in the next as a general
algebraic operation on vectors, continuing this book's movement from
individual geometric configurations toward increasingly unified
descriptions of transformation and space.

\section{Exercises}

\subsection*{The Law of Sines}
\begin{exercise}
For each triangle, use the Law of Sines to find the missing side, rounded
to the nearest tenth: (a) $A=35^\circ$, $B=85^\circ$, $a=12$; (b)
$A=48^\circ$, $C=67^\circ$, $c=15$; (c) $B=72^\circ$, $C=41^\circ$,
$b=18$.
\end{exercise}

\begin{exercise}
A triangle has $A = 40^\circ$, $B = 65^\circ$, $a = 10$. Find $C$, $b$,
and $c$.
\end{exercise}

\begin{exercise}
A triangle has $B = 58^\circ$, $C = 47^\circ$, $c = 14$. Solve the
triangle completely.
\end{exercise}

\begin{exercise}
Verify directly that the triangle with $A = 30^\circ$, $B = 60^\circ$, $C
= 90^\circ$ satisfies the Law of Sines.
\end{exercise}

\subsection*{The Ambiguous Case}
\begin{exercise}
Determine whether each SSA configuration produces no triangle, one
triangle, or two triangles: (a) $A=35^\circ$, $a=8$, $b=15$; (b)
$A=40^\circ$, $a=12$, $b=10$; (c) $A=25^\circ$, $a=9$, $b=18$.
\end{exercise}

\begin{exercise}
Solve completely, giving all possible triangles: $A = 30^\circ$, $a = 10$,
$b = 18$.
\end{exercise}

\begin{exercise}
Solve completely: $A = 50^\circ$, $a = 16$, $b = 12$.
\end{exercise}

\begin{exercise}
Explain why the SSA case can sometimes produce two different triangles.
\end{exercise}

\subsection*{The Law of Cosines}
\begin{exercise}
Use the Law of Cosines to find the missing side: (a) $a=8$, $b=11$,
$C=50^\circ$; (b) $a=13$, $b=9$, $C=120^\circ$; (c) $a=7$, $b=14$,
$C=35^\circ$.
\end{exercise}

\begin{exercise}
Find the third side of a triangle whose sides of length $12$ and $17$
include an angle of $110^\circ$.
\end{exercise}

\begin{exercise}
Two roads leave a town and diverge at an angle of $75^\circ$. A traveler
goes $18$ km along one road and another traveler goes $25$ km along the
other. How far apart are they?
\end{exercise}

\subsection*{Finding Angles from Three Sides}
\begin{exercise}
Use the Law of Cosines to find the indicated angle: (a) $a=7$, $b=9$,
$c=11$, find $C$; (b) $a=10$, $b=10$, $c=12$, find $C$; (c) $a=5$, $b=8$,
$c=10$, find $A$.
\end{exercise}

\begin{exercise}
Determine all three angles of a triangle with sides $6$, $8$, $10$.
\end{exercise}

\begin{exercise}
Determine all three angles of a triangle with sides $9$, $11$, $15$.
\end{exercise}

\begin{exercise}
Explain why the SSS case always determines a unique triangle whenever the
triangle inequalities are satisfied.
\end{exercise}

\subsection*{Area of a Triangle}
\begin{exercise}
Find the area: (a) $a=12$, $b=15$, $C=40^\circ$; (b) $a=9$, $b=14$,
$C=65^\circ$; (c) $b=20$, $c=11$, $A=32^\circ$.
\end{exercise}

\begin{exercise}
A triangular plot of land has sides of length $60$ m and $85$ m with an
included angle of $54^\circ$. Find its area.
\end{exercise}

\begin{exercise}
Show that $\tfrac12 ab\sin C = \tfrac12 bc\sin A$ using the Law of Sines.
\end{exercise}

\subsection*{Mixed Problems}
\begin{exercise}
Solve each triangle completely: (a) $a=8$, $b=11$, $C=50^\circ$; (b)
$A=42^\circ$, $B=68^\circ$, $a=12$; (c) $a=7$, $b=9$, $c=11$.
\end{exercise}

\begin{exercise}
A surveyor stands at point $P$ and measures the angle between two
landmarks $Q$ and $R$ to be $63^\circ$. The distances from $P$ to the
landmarks are $240$ m and $310$ m. Find the distance between the
landmarks.
\end{exercise}

\begin{exercise}
A ship travels $120$ km due east and then changes course, traveling $90$
km on a heading $35^\circ$ north of east. How far is the ship from its
starting point?
\end{exercise}

\begin{exercise}
A triangular support frame has side lengths $4.2$ m, $5.6$ m, and $7.1$
m. Determine its largest angle.
\end{exercise}

\begin{exercise}
Explain when the Law of Sines is usually preferable to the Law of
Cosines, and when the reverse is true.
\end{exercise}

\subsection*{Conceptual and Exploratory Problems}
\begin{exercise}
Show algebraically that the Law of Cosines reduces to the Pythagorean
theorem when $C = 90^\circ$.
\end{exercise}

\begin{exercise}
Describe geometrically the meaning of the correction term $-2ab\cos C$ in
the Law of Cosines.
\end{exercise}

\begin{exercise}
A student claims that the Law of Cosines is merely a complicated version
of the Pythagorean theorem. Evaluate this claim.
\end{exercise}

\begin{exercise}
The Law of Sines and the Law of Cosines both reconstruct unknown
measurements from partial information. Explain how each accomplishes this
in a different way.
\end{exercise}

\begin{exercise}
Describe how this chapter extends the trigonometry of right triangles to
arbitrary triangles.
\end{exercise}

\subsection*{Challenge Problems}
\begin{exercise}
A triangle has sides $13$, $14$, $15$. Find its area without using
Heron's formula.
\end{exercise}

\begin{exercise}
A triangle has $A = 25^\circ$, $a = 8$, $b = 14$. Determine all possible
triangles and classify each as acute, right, or obtuse.
\end{exercise}

\begin{exercise}
Prove the Law of Cosines using an altitude rather than coordinates.
\end{exercise}

\begin{exercise}
For a triangle with sides $a$, $b$, $c$, show that
$\cos C = \dfrac{a^2+b^2-c^2}{2ab}$, and explain why this formula allows
any angle of a triangle to be computed directly from the side lengths
alone.
\end{exercise}
